% !TEX root = 第四章_线性空间与线性方程组(答案版).tex
\setcounter{chapter}{3}
\pagestyle{fancy}

\chapter{线性空间与线性方程组}

\section{线性空间与向量组的线性相关性}

%例题4.1.1
\begin{example}[\markstar]
若向量组 $\alpha_1,\alpha_2,\ldots,\alpha_m$ 线性无关,而
$\beta,\alpha_1,\alpha_2,\ldots,\alpha_m$ 线性相关,证明:$\beta$ 可由
$\alpha_1,\ldots,\alpha_m$ 线性表示,且表示式唯一.
\end{example}
\exampleblank

\begin{solution}
由线性相关性,存在不全为零的数 $c,c_1,\ldots,c_m$ 使
\[
 c\beta+\sum\limits_{i=1}^m c_i\alpha_i=0.
\]
若 $c=0$,则 $\displaystyle\sum c_i\alpha_i=0$ 与 $\alpha_1,\ldots,\alpha_m$ 线性无关矛盾,故 $c\neq0$,从而
\[
 \beta=-\sum\limits_{i=1}^m\dfrac{c_i}{c}\alpha_i.
\]
若还有另一表示,两式相减即得到 $\alpha_1,\ldots,\alpha_m$ 的非平凡线性关系,故表示唯一.
\end{solution}

%例题4.1.2
\begin{example}[\markstar]
证明:若 $\alpha_1,\ldots,\alpha_r$ 线性相关,则
$\alpha_1,\ldots,\alpha_r,\alpha_{r+1},\ldots,\alpha_n$ 也线性相关.
\end{example}
\exampleblank

\begin{solution}
在 $\alpha_1,\ldots,\alpha_r$ 的非平凡线性关系后补零系数即可得到扩大向量组的非平凡线性关系.
\end{solution}

%例题4.1.3
\begin{example}[\markstar]
证明:若一组 $r$ 维向量线性无关,则给每个向量添加相同位置的 $n-r$ 个分量后所得的 $n$ 维向量组仍线性无关.
\end{example}
\exampleblank

\begin{solution}
设扩充后的向量为 $\widetilde\alpha_i$.若
\[
 \displaystyle\sum_i c_i\widetilde\alpha_i=0,
\]
只看原来的前 $r$ 个分量,就有 $\displaystyle\sum_i c_i\alpha_i=0$.原向量组线性无关,故所有 $c_i=0$.
\end{solution}

%例题4.1.4
\begin{example}
设 $\alpha_1,\ldots,\alpha_m$ 线性无关,$\beta_1$ 可由它们线性表示,而 $\beta_2$ 不能由它们线性表示.证明:对任意常数 $l$,
\[
 \alpha_1,\ldots,\alpha_m,l\beta_1+\beta_2
\]
总线性无关.
\end{example}
\exampleblank

\begin{solution}
若该组线性相关,则因 $\alpha_1,\ldots,\alpha_m$ 线性无关,最后一个向量必可由它们线性表示.又 $l\beta_1$ 已可由它们表示,因此 $\beta_2$ 也可由它们表示,矛盾.
\end{solution}

%例题4.1.5
\begin{example}
设 $\alpha_1,\ldots,\alpha_m$ 线性无关,而 $\alpha_2,\ldots,\alpha_{m+1}$ 线性相关.证明:$\alpha_1$ 不能由 $\alpha_2,\ldots,\alpha_{m+1}$ 线性表出.
\end{example}
\exampleblank

\begin{solution}
若 $\alpha_1$ 可由 $\alpha_2,\ldots,\alpha_{m+1}$ 表出,而 $\alpha_2,\ldots,\alpha_{m+1}$ 线性相关,则去掉其中一个可由其余向量表示的向量,可推出 $\alpha_1,\ldots,\alpha_m$ 线性相关,矛盾.等价地,也可将两条线性关系消元得到 $\alpha_1,\ldots,\alpha_m$ 的非平凡关系.
\end{solution}

%例题4.1.6
\begin{example}
设 $\alpha_1,\ldots,\alpha_s$ 线性无关.判断
\[
 \alpha_1+\alpha_2,\ \alpha_2+\alpha_3,\ \ldots,\ \alpha_{s-1}+\alpha_s,\ \alpha_s+\alpha_1
\]
是否线性无关.
\end{example}
\exampleblank

\begin{solution}
设这些向量的系数为 $c_1,\ldots,c_s$.由原向量组线性无关,得到
\[
 c_1+c_s=0,\qquad c_1+c_2=0,\qquad\ldots,\qquad c_{s-1}+c_s=0.
\]
于是 $c_{i+1}=-c_i$,并有 $c_s=(-1)^{s-1}c_1$ 与 $c_s=-c_1$.故当 $s$ 为奇数时只能 $c_1=0$,向量组线性无关;当 $s$ 为偶数时可取
\[
 (c_1,c_2,\ldots,c_s)=(1,-1,1,-1,\ldots,-1),
\]
得到非平凡关系.因此
\[
 \boxed{\text{线性无关当且仅当 }s\text{ 为奇数}.}
\]
\end{solution}

%例题4.1.7
\begin{example}[\markstar]
设 $\alpha_1,\ldots,\alpha_r$ 线性无关,且
\[
 \beta_i=\sum\limits_{j=1}^r k_{ij}\alpha_j,\qquad i=1,\ldots,r.
\]
证明:$\beta_1,\ldots,\beta_r$ 线性相关的充要条件是系数矩阵 $K=(k_{ij})$ 不可逆.
\end{example}
\exampleblank

\begin{solution}
若 $c=(c_1,\ldots,c_r)^{\mathsf T}$,则
\[
 \displaystyle\sum_i c_i\beta_i
 =\sum_j (K^{\mathsf T}c)_j\alpha_j.
\]
因 $\alpha_j$ 线性无关,左端为零当且仅当 $K^{\mathsf T}c=0$.存在非零 $c$ 当且仅当 $K$ 不可逆.
\end{solution}

%例题4.1.8
\begin{example}[\markstar]
设 $\alpha_1,\ldots,\alpha_m$ 线性无关,
\[
 \beta_i=\sum\limits_{j=1}^m a_{ij}\alpha_j,\qquad i=1,\ldots,k,
\]
记 $A=(a_{ij})_{k\times m}$.证明:向量组 $\beta_1,\ldots,\beta_k$ 的秩等于 $r(A)$.
\end{example}
\exampleblank

\begin{solution}
映射
\[
 (c_1,\ldots,c_m)\longmapsto \sum\limits_{j=1}^m c_j\alpha_j
\]
因 $\alpha_j$ 线性无关而为单射.$\beta_i$ 的坐标行正是 $A$ 的各行,因此这些向量之间的线性关系与 $A$ 的行向量之间的线性关系完全一致,故秩相等.
\end{solution}

%例题4.1.9
\begin{example}
设 $\alpha_1,\ldots,\alpha_s,\beta$ 为 $s+1$ 个向量,$s>1$,且
\[
 \beta=\alpha_1+\cdots+\alpha_s.
\]
证明:$\beta-\alpha_1,\ldots,\beta-\alpha_s$ 线性无关的充要条件是 $\alpha_1,\ldots,\alpha_s$ 线性无关.
\end{example}
\exampleblank

\begin{solution}
记 $\gamma_i=\beta-\alpha_i$.由
\[
 \displaystyle\sum\limits_{i=1}^s\gamma_i=(s-1)\beta
\]
得
\[
 \beta=\dfrac{1}{s-1}\sum_i\gamma_i,
 \qquad
 \alpha_i=\beta-\gamma_i.
\]
故两组向量可以互相线性表示.它们的向量个数相同,因此一组线性无关当且仅当另一组线性无关.
\end{solution}

\begin{theorem}[\markstar]
对于向量组 $A:\alpha_1,\ldots,\alpha_r$ 与 $B:\beta_1,\ldots,\beta_s$,若 $A$ 线性无关且 $A$ 可由 $B$ 线性表示,则 $r\leq s$.
\end{theorem}

\begin{corollary}[\markstar]
若 $m>n$,则任意 $m$ 个 $n$ 维向量必线性相关.
\end{corollary}

\begin{theorem}[\markstar]
若向量组 $A$ 的秩为 $r_1$,向量组 $B$ 的秩为 $r_2$,且 $A$ 可由 $B$ 线性表示,则 $r_1\leq r_2$.
\end{theorem}

\begin{corollary}[\markstar]
等价向量组具有相同的秩;任一向量组的任意两个极大无关组所含向量个数相同.
\end{corollary}

%例题4.1.10
\begin{example}[\markstar]
设 $\alpha_1,\ldots,\alpha_s$ 线性无关,并且可由 $\beta_1,\ldots,\beta_t$ 线性表示.证明:必存在某个 $\beta_j$,使
\[
 \beta_j,\alpha_2,\ldots,\alpha_s
\]
线性无关.
\end{example}
\exampleblank

\begin{solution}
若对所有 $j$,$\beta_j,\alpha_2,\ldots,\alpha_s$ 都线性相关,则每个 $\beta_j$ 都可由 $\alpha_2,\ldots,\alpha_s$ 线性表示.于是由 $\beta_j$ 可表示的 $\alpha_1$ 也可由 $\alpha_2,\ldots,\alpha_s$ 表示,矛盾.
\end{solution}

%例题4.1.11
\begin{example}
设 $A=(\alpha_1,\ldots,\alpha_r)$ 为 $n\times r$ 矩阵,$B=(\beta_1,\ldots,\beta_s)$ 为 $n\times s$ 矩阵,且 $r(A)=r,r(B)=s$.证明:若 $r+s>n$,则存在非零向量 $\xi$,既可由 $\alpha_i$ 线性表示,也可由 $\beta_j$ 线性表示.
\end{example}
\exampleblank

\begin{solution}
令 $U=L(\alpha_1,\ldots,\alpha_r)$,$W=L(\beta_1,\ldots,\beta_s)$.由维数公式
\[
 \dim(U\cap W)=\dim U+\dim W-\dim(U+W)
 \ge r+s-n>0.
\]
故 $U\cap W$ 中存在非零向量 $\xi$.
\end{solution}

%例题4.1.12
\begin{example}
设
\[
V=\{(a+ib,c+id)\mid a,b,c,d\in\mathbb{R}\}.
\]
分别把它看作复数域和实数域上的线性空间 $V_{\mathbb{C}},V_{\mathbb{R}}$,求其维数.又设 $M_2(\mathbb{C})$ 分别按复数域和实数域看作线性空间 $U_{\mathbb{C}},U_{\mathbb{R}}$,求其维数.
\end{example}
\exampleblank

\begin{solution}
显然 $V\cong\mathbb{C}^2$,故
\[
 \dim_{\mathbb{C}}V_{\mathbb{C}}=2,
 \qquad
 \dim_{\mathbb{R}}V_{\mathbb{R}}=4.
\]
同理 $M_2(\mathbb{C})\cong\mathbb{C}^4$,故
\[
 \dim_{\mathbb{C}}U_{\mathbb{C}}=4,
 \qquad
 \dim_{\mathbb{R}}U_{\mathbb{R}}=8.
\]
\end{solution}

%例题4.1.13
\begin{example}
设 $V_1,V_2$ 是数域 $P$ 上的线性空间,在 $V_1\times V_2$ 上按分量定义加法与数乘.证明它构成线性空间;若 $\dim V_1=m,\dim V_2=n$,求 $\dim(V_1\times V_2)$.
\end{example}
\exampleblank

\begin{solution}
线性空间公理由两个分量逐项验证.若 $\alpha_1,\ldots,\alpha_m$ 与 $\beta_1,\ldots,\beta_n$ 分别为 $V_1,V_2$ 的基,则
\[
 (\alpha_1,0),\ldots,(\alpha_m,0),(0,\beta_1),\ldots,(0,\beta_n)
\]
是 $V_1\times V_2$ 的一组基.因此
\[
 \boxed{\dim(V_1\times V_2)=m+n}.
\]
\end{solution}

%例题4.1.14
\begin{example}
在二维实向量集合 $V$ 上规定
\[
(a,b)\oplus(c,d)=(a+c,b+d+ac),
\]
\[
k\circ(a,b)=\left(ka,kb+\dfrac{k(k-1)}{2}a^2\right).
\]
证明这构成 $\mathbb{R}$ 上的线性空间,并讨论 $\alpha_1=(1,1),\alpha_2=(a,b)$ 的线性相关性.
\end{example}
\exampleblank

\begin{solution}
定义
\[
 \Phi(a,b)=\left(a,b-\dfrac{a^2}{2}\right).
\]
直接计算得
\[
 \Phi(u\oplus v)=\Phi(u)+\Phi(v),
 \qquad
 \Phi(k\circ u)=k\Phi(u),
\]
故 $V$ 与通常的 $\mathbb{R}^2$ 线性同构,从而确为线性空间.

又
\[
 \Phi(\alpha_1)=\left(1,\dfrac12\right),
 \qquad
 \Phi(\alpha_2)=\left(a,b-\dfrac{a^2}{2}\right).
\]
两向量线性相关当且仅当
\[
 b-\dfrac{a^2}{2}-\dfrac a2=0,
\]
即
\[
 \boxed{2b=a^2+a}.
\]
\end{solution}

%例题4.1.15
\begin{example}
已知
\[
\alpha_1=(1,2,4,3)^{\mathsf T},\quad
\alpha_2=(1,-1,-6,6)^{\mathsf T},
\]
\[
\alpha_3=(-2,-1,2,-9)^{\mathsf T},\quad
\alpha_4=(1,1,-2,7)^{\mathsf T},
\]
\[
\beta=(4,2,4,a)^{\mathsf T}.
\]
设 $W=L(\alpha_1,\alpha_2,\alpha_3,\alpha_4)$.
\begin{enumerate}
\item 求 $\dim W$ 及一组基;
\item 求 $a$ 使 $\beta\in W$,并求 $\beta$ 在所选基下的坐标.
\end{enumerate}
\end{example}
\exampleblank

\begin{solution}
对列矩阵作行化简可得主元列为第 $1,2,4$ 列,故
\[
 \dim W=3,
 \qquad
 \{\alpha_1,\alpha_2,\alpha_4\}
\]
可取为一组基.

计算
\[
 \det(\alpha_1,\alpha_2,\alpha_4,\beta)=8(a-9),
\]
故 $\beta\in W$ 当且仅当 $a=9$.此时
\[
 \beta=4\alpha_1+3\alpha_2-3\alpha_4,
\]
所以坐标为
\[
 \boxed{(4,3,-3)^{\mathsf T}}.
\]
\end{solution}

%例题4.1.16
\begin{example}
设
\[
W=\{f(x)\in\mathbb{R}[x]_n\mid f(1)=0\},
\]
其中 $\mathbb{R}[x]_n$ 表示次数小于 $n$ 的多项式及零多项式.证明 $W$ 是子空间,并求其维数和一组基.
\end{example}
\exampleblank

\begin{solution}
条件 $f(1)=0$ 对加法和数乘封闭,故 $W$ 为子空间.由因式定理,
\[
 f(1)=0\iff f(x)=(x-1)g(x),\qquad \deg g<n-1.
\]
故
\[
 W=(x-1)\mathbb{R}[x]_{n-1},
 \qquad \dim W=n-1,
\]
可取基
\[
 \boxed{x-1,x(x-1),\ldots,x^{n-2}(x-1)}.
\]
\end{solution}

%例题4.1.17
\begin{example}
设 $M_n(K)$ 为数域 $K$ 上的 $n$ 阶方阵空间,$U$ 为全体循环矩阵组成的集合.证明 $U$ 是 $M_n(K)$ 的子空间,并求一组基和维数.
\end{example}
\exampleblank

\begin{solution}
设 $P$ 为循环移位矩阵,则任一循环矩阵唯一写成
\[
 a_0E+a_1P+\cdots+a_{n-1}P^{n-1}.
\]
因此 $U$ 对加法和数乘封闭,且
\[
 \boxed{E,P,\ldots,P^{n-1}}
\]
为一组基,从而
\[
 \boxed{\dim U=n}.
\]
\end{solution}

%例题4.1.18
\begin{example}
设 $A=J_n(\lambda)$ 为一个 $n$ 阶 Jordan 块,
\[
V=\{B\in M_n(\mathbb{R})\mid BA=AB\}.
\]
证明 $V$ 是线性空间,并求 $\dim V$.
\end{example}
\exampleblank

\begin{solution}
可交换条件对加法和数乘封闭,故 $V$ 为子空间.写
\[
 A=\lambda E+N,
\]
其中 $N$ 为幂零 Jordan 移位矩阵.$BA=AB$ 等价于 $BN=NB$.比较矩阵元素可知 $B$ 必为上三角 Toeplitz 矩阵,即
\[
 B=c_0E+c_1N+\cdots+c_{n-1}N^{n-1}.
\]
故
\[
 \boxed{\dim V=n}.
\]
\end{solution}


\newpage

\section{向量组与矩阵的秩}

\begin{proposition}[\markstar\ 求极大无关组]
把向量组作为列向量排成矩阵,对矩阵作初等行变换化为阶梯形;主元列在原矩阵中对应的向量可组成一个极大无关组.
\end{proposition}

\begin{remark}
初等行变换不改变列向量之间的线性关系,因此也可据此求向量的线性表示坐标.极大无关组通常不唯一.
\end{remark}

%例题4.2.1
\begin{example}
求下列向量组的秩,并给出一个极大无关组:
\begin{enumerate}
  \item $(1,2,3),(4,8,12),(3,0,1),(4,5,8)$;
  \item $(1,2,3,4),(1,0,1,0),(-1,1,1,-1),(-2,3,2,3)$;
  \item $(4,-5,2,6),(2,-2,1,3),(6,-3,3,9),(4,-1,5,6)$;
  \item $(1,0,0,2,5),(0,1,0,3,4),(0,0,1,4,7),(2,-3,4,11,12)$.
\end{enumerate}
\end{example}
\exampleblank

\begin{solution}
把每组向量作为列排成矩阵并作行化简.可分别取下列主元列:
\begin{enumerate}
  \item 第 $1,3$ 列,故秩为 $2$;
  \item 第 $1,2,3$ 列,故秩为 $3$;
  \item 第 $1,2,4$ 列,故秩为 $3$;
  \item 四列全为主元列,故秩为 $4$.
\end{enumerate}
对应原向量即为所求极大无关组.
\end{solution}

%例题4.2.2
\begin{example}
设向量组
\[
(I):\quad
\alpha_1=(1,1,a)^{\mathsf T},\quad
\alpha_2=(-2,a,4)^{\mathsf T},\quad
\alpha_3=(-2,a,a)^{\mathsf T},
\]
\[
(II):\quad
\beta_1=(1,1,a)^{\mathsf T},\quad
\beta_2=(1,a,1)^{\mathsf T},\quad
\beta_3=(a,1,1)^{\mathsf T}.
\]
\begin{enumerate}
  \item 求 $a$ 使向量组 $(I)$ 线性相关;
  \item 求 $a$ 使向量组 $(I)$ 不能由向量组 $(II)$ 线性表示;
  \item 在前两问同时成立时,将 $\gamma=(1,-2,-5)^{\mathsf T}$ 用 $\beta_1,\beta_2,\alpha_3$ 线性表示.
\end{enumerate}
\end{example}
\exampleblank

\begin{solution}
有
\[
 \det(\alpha_1,\alpha_2,\alpha_3)=(a-4)(a+2),
\]
故第一问为
\[
 a=4\quad\text{或}\quad a=-2.
\]
又
\[
 \det(\beta_1,\beta_2,\beta_3)=-(a-1)^2(a+2).
\]
当 $a=4$ 时,$(II)$ 为一组基,故可表示任意三维向量;于是前两问同时成立只能取
\[
 \boxed{a=-2}.
\]
此时直接解线性方程得
\[
 \boxed{\gamma=2\beta_1+\beta_2+\alpha_3}.
\]
\end{solution}

\begin{proposition}[\marktwostar]
若 $A_{m\times n},B_{n\times s}$,则
\[
 r(AB)\leq\min\{r(A),r(B)\}.
\]
\end{proposition}

\begin{proposition}[\marktwostar]
若 $A,B$ 同型,则
\[
 r(A\pm B)\leq r(A)+r(B).
\]
\end{proposition}

\begin{proposition}[\marktwostar]
对分块对角矩阵有
\[
 r\begin{pmatrix}A&O\\O&B\end{pmatrix}=r(A)+r(B).
\]
\end{proposition}

\begin{proposition}[\marktwostar]
对适当阶数的分块矩阵,
\[
 r\begin{pmatrix}A&O\\C&B\end{pmatrix}\geq r(A)+r(B).
\]
\end{proposition}

\begin{proposition}[\markstar]
若 $A_{m\times n}B_{n\times s}=O$,则
\[
 r(A)+r(B)\leq n.
\]
\end{proposition}

\begin{lemma}[\marktwostar]
若 $\alpha$ 是 $n$ 维实列向量,则
\[
 \alpha^{\mathsf T}\alpha=0\iff \alpha=0.
\]
复向量情形必须把普通转置改为共轭转置.
\end{lemma}

\begin{proposition}[\marktwostar]
若 $A$ 为实矩阵,则
\[
 r(A^{\mathsf T}A)=r(AA^{\mathsf T})=r(A).
\]
\end{proposition}

%例题4.2.3
\begin{example}
设 $A_{m\times n},B_{s\times n}$,证明
\[
 r\begin{pmatrix}A\\B\end{pmatrix}\leq r(A)+r(B).
\]
\end{example}
\exampleblank

\begin{solution}
矩阵的行空间包含于 $A$ 的行空间与 $B$ 的行空间之和,故
\[
 r\begin{pmatrix}A\\B\end{pmatrix}
 =\dim(R(A)+R(B))
 \leq r(A)+r(B).
\]
\end{solution}

%例题4.2.4
\begin{example}
设 $A_{m\times n},B_{m\times s}$,证明
\[
 r(A,B)\leq r(A)+r(B).
\]
\end{example}
\exampleblank

\begin{solution}
同理,$(A,B)$ 的列空间为 $C(A)+C(B)$,故
\[
 r(A,B)=\dim(C(A)+C(B))\leq r(A)+r(B).
\]
\end{solution}

%例题4.2.5
\begin{example}
设 $A,B$ 为 $n$ 阶幂等矩阵,且 $E-A-B$ 可逆.证明
\[
 r(A)=r(B).
\]
\end{example}
\exampleblank

\begin{solution}
令 $T=E-A-B$.若 $x\in\operatorname{Im}A$,则 $Ax=x$,故
\[
 Tx=-Bx\in\operatorname{Im}B.
\]
因为 $T$ 可逆,$T$ 在 $\operatorname{Im}A$ 上的限制是单射,所以
\[
 r(A)\leq r(B).
\]
交换 $A,B$ 得 $r(B)\leq r(A)$,于是
\[
 \boxed{r(A)=r(B)}.
\]
\end{solution}

%例题4.2.6
\begin{example}
设 $A,B$ 为 $n$ 阶方阵,证明
\[
 r(AB-E)\leq r(A-E)+r(B-E).
\]
\end{example}
\exampleblank

\begin{solution}
由
\[
 AB-E=A(B-E)+(A-E)
\]
及秩的不等式,
\[
 r(AB-E)
 \leq r(A(B-E))+r(A-E)
 \leq r(B-E)+r(A-E).
\]
\end{solution}

%例题4.2.7
\begin{example}
设 $A,B$ 为 $n$ 阶方阵,证明:
\begin{enumerate}
  \item $r(A-B)\geq r(A)-r(B)$;
  \item 若 $A$ 可逆,则第一问等号成立当且仅当
  \[
  BA^{-1}B=B.
  \]
\end{enumerate}
\end{example}
\exampleblank

\begin{solution}
由 $A=(A-B)+B$ 得
\[
 r(A)\leq r(A-B)+r(B),
\]
即第一式.

当 $A$ 可逆时,令 $C=BA^{-1}$,则
\[
 r(A-B)=r((E-C)A)=r(E-C),\qquad r(B)=r(C).
\]
总有
\[
 \ker(E-C)\subseteq\operatorname{Im}C.
\]
因此
\[
 r(E-C)=n-r(C)
\]
当且仅当 $\ker(E-C)=\operatorname{Im}C$,这又等价于 $C^2=C$.故
\[
 C^2=C
 \iff BA^{-1}BA^{-1}=BA^{-1}
 \iff \boxed{BA^{-1}B=B}.
\]
\end{solution}

%例题4.2.8
\begin{example}
设 $A,D$ 分别为 $m,n$ 阶可逆矩阵,$B,C$ 分别为 $m\times n,n\times m$ 矩阵.证明
\[
 r(A)-r(A-BD^{-1}C)
 =r(D)-r(D-CA^{-1}B).
\]
\end{example}
\exampleblank

\begin{solution}
对分块矩阵
\[
 M=\begin{pmatrix}A&B\\C&D\end{pmatrix}
\]
分别以 $A,D$ 作 Schur 消元,有
\[
 r(M)=m+r(D-CA^{-1}B)
\]
和
\[
 r(M)=n+r(A-BD^{-1}C).
\]
两式相等并整理即得结论.
\end{solution}

\begin{corollary}[\markstar]
特别地,令 $A=\lambda_0E_m,D=E_n$,则
\[
 m-r(\lambda_0E_m-BC)
 =n-r(\lambda_0E_n-CB).
\]
\end{corollary}

%例题4.2.9
\begin{example}[\markstar]
设 $A$ 为 $m\times n$ 矩阵,$B$ 为 $n\times s$ 矩阵,且
\[
 r(AB)=r(B).
\]
证明:对任意 $s\times t$ 矩阵 $C$,有
\[
 r(ABC)=r(BC).
\]
\end{example}
\exampleblank

\begin{solution}
条件 $r(AB)=r(B)$ 等价于线性映射 $A$ 在 $\operatorname{Im}B$ 上为单射.由于
\[
 \operatorname{Im}(BC)\subseteq\operatorname{Im}B,
\]
$A$ 在 $\operatorname{Im}(BC)$ 上仍单射,故
\[
 \dim A(\operatorname{Im}(BC))=\dim\operatorname{Im}(BC),
\]
即 $r(ABC)=r(BC)$.
\end{solution}

\begin{proposition}[\marktwostar]
设 $A$ 是 $n$ 阶方阵,则存在 $1\leq k\leq n$,使
\[
 r(A^k)=r(A^{k+1})=r(A^{k+2})=\cdots.
\]
\end{proposition}

%例题4.2.10
\begin{example}
设 $A$ 为 $n$ 阶实方阵,满足
\[
 AA^{\mathsf T}=a^2E_n.
\]
证明
\[
 r(aE_n-A)=r((aE_n-A)^2).
\]
\end{example}
\exampleblank

\begin{solution}
若 $a=0$,则 $A=0$,结论显然.以下设 $a\neq0$,令 $Q=A/a$,则 $QQ^{\mathsf T}=E$,即 $Q$ 为正交矩阵.

只需证明
\[
 \ker(E-Q)^2=\ker(E-Q).
\]
若 $(E-Q)^2x=0$,令 $y=(E-Q)x$,则 $Qy=y$.另一方面,正交矩阵满足
\[
 \operatorname{Im}(E-Q)\perp\ker(E-Q),
\]
而 $y\in\operatorname{Im}(E-Q)\cap\ker(E-Q)$,故 $y=0$.于是 $(E-Q)x=0$,两核相同,因此两矩阵秩相同.
\end{solution}

%例题4.2.11
\begin{example}
\begin{enumerate}
  \item 设 $A,B$ 都是 $n$ 阶方阵且 $AB=O$,证明
  \[
  r(BA)\leq\left\lfloor\dfrac n2\right\rfloor;
  \]
  \item 对任意正整数 $n$,构造 $n$ 阶矩阵 $A,B$,使 $AB=O$ 且
  \[
  r(BA)=\left\lfloor\dfrac n2\right\rfloor.
  \]
\end{enumerate}
\end{example}
\exampleblank

\begin{solution}
由 $AB=O$ 得
\[
 (BA)^2=B(AB)A=O.
\]
所以
\[
 \operatorname{Im}(BA)\subseteq\ker(BA),
\]
从而
\[
 r(BA)\leq n-r(BA),
\]
即第一式.

令 $k=\lfloor n/2\rfloor$,按分解 $\mathbb{F}^n=U\oplus W$,其中 $\dim U=k$.取
\[
 B=\begin{pmatrix}E_k&O\\O&O\end{pmatrix},
 \qquad
 A=\begin{pmatrix}O&C\\O&O\end{pmatrix},
\]
其中 $C$ 为 $k\times(n-k)$ 的满行秩矩阵.则 $AB=O$,而 $BA=A$,故
\[
 r(BA)=r(A)=k.
\]
\end{solution}

%例题4.2.12
\begin{example}[\markstar\ Sylvester 不等式]
设 $A_{m\times n},B_{n\times s}$,证明
\[
 r(AB)\geq r(A)+r(B)-n.
\]
\end{example}
\exampleblank

\begin{solution}
把 $A$ 限制在 $\operatorname{Im}B$ 上.由秩--零度定理,
\[
 r(AB)
 =\dim\operatorname{Im}B-
 \dim(\operatorname{Im}B\cap\ker A).
\]
而
\[
 \dim(\operatorname{Im}B\cap\ker A)
 \leq\dim\ker A=n-r(A).
\]
故
\[
 r(AB)\geq r(B)-(n-r(A)).
\]
\end{solution}

%例题4.2.13
\begin{example}[\markstar\ Frobenius 不等式]
设 $A_{m\times n},B_{n\times s},C_{s\times t}$,证明
\[
 r(ABC)\geq r(AB)+r(BC)-r(B).
\]
\end{example}
\exampleblank

\begin{solution}
把 $A$ 限制在 $\operatorname{Im}B$ 上,并进一步限制到 $\operatorname{Im}(BC)$.有
\[
 r(ABC)=\dim\operatorname{Im}(BC)-\dim(\operatorname{Im}(BC)\cap\ker A),
\]
而
\[
 \operatorname{Im}(BC)\cap\ker A
 \subseteq \operatorname{Im}B\cap\ker A.
\]
又
\[
 \dim(\operatorname{Im}B\cap\ker A)=r(B)-r(AB).
\]
故
\[
 r(ABC)\geq r(BC)-r(B)+r(AB).
\]
\end{solution}

%例题4.2.14
\begin{example}
设 $A$ 为 $n$ 阶方阵,证明:
\begin{enumerate}
  \item $A^2=E$ 当且仅当
  \[
  r(A+E)+r(A-E)=n;
  \]
  \item $A^2=A$ 当且仅当
  \[
  r(A)+r(E-A)=n.
  \]
\end{enumerate}
\end{example}
\exampleblank

\begin{solution}
第一问中总有
\[
 n=r(2E)=r((A+E)-(A-E))\leq r(A+E)+r(A-E).
\]
若 $A^2=E$,则 $(A+E)(A-E)=O$,由秩不等式得反向不等式,因此等号成立.

反之若秩和等于 $n$,则
\[
 \operatorname{Im}(A+E)\cap\operatorname{Im}(A-E)=\{0\}.
\]
而 $(A+E)(A-E)$ 的像同时包含于这两个像空间,故只能为零,于是 $A^2=E$.

第二问同理.由
\[
 E=A+(E-A)
\]
知秩和至少为 $n$;若 $A(E-A)=O$ 则至多为 $n$.反过来若秩和恰为 $n$,两像空间交为零,而 $A(E-A)$ 的像属于二者之交,故 $A(E-A)=O$,即 $A^2=A$.
\end{solution}

%例题4.2.15
\begin{example}
设 $A^3=E$,证明
\[
 r(A-E)+r(A^2+A+E)=n.
\]
进一步,若 $A\in M_n(\mathbb{C})$,$f,g\in\mathbb{C}[x]$ 且 $(f,g)=1$,证明
\[
 n=r(f(A))+r(g(A))-r(f(A)g(A)).
\]
\end{example}
\exampleblank

\begin{solution}
先证推广式.由 Bézout 恒等式,存在 $u,v\in\mathbb{C}[x]$ 使
\[
 u(A)f(A)+v(A)g(A)=E.
\]
于是
\[
 \operatorname{Im}f(A)+\operatorname{Im}g(A)=\mathbb{C}^n.
\]
又因为 $f(A),g(A)$ 可交换,若
\[
 y\in\operatorname{Im}f(A)\cap\operatorname{Im}g(A),
\]
写 $y=f(A)x=g(A)z$,则
\[
 y=u(A)f(A)g(A)z+v(A)g(A)f(A)x
 \in\operatorname{Im}(f(A)g(A)).
\]
反向包含显然,故
\[
 \operatorname{Im}f(A)\cap\operatorname{Im}g(A)
 =\operatorname{Im}(f(A)g(A)).
\]
由维数公式即得推广式.

取
\[
 f(x)=x-1,\qquad g(x)=x^2+x+1.
\]
在特征为零的数域上二者互素,且 $f(A)g(A)=A^3-E=O$,故得到所求秩和等式.
\end{solution}

\subsection{满秩矩阵}

\begin{definition}
若 $C_{m\times n}$ 满足 $r(C)=m$,称为行满秩矩阵;若 $r(C)=n$,称为列满秩矩阵.
\end{definition}

\begin{theorem}[\marktwostar]
设 $C$ 为列满秩矩阵,则下列命题等价:
\begin{enumerate}
  \item $Cx=0$ 只有零解;
  \item $C$ 保持线性无关组的线性无关性;
  \item $C$ 有左逆;
  \item $C$ 可扩充成可逆矩阵;
  \item 对任意适当维数的 $b$,$C^{\mathsf T}x=b$ 总有解.
\end{enumerate}
转置后得到行满秩矩阵的对应结论.
\end{theorem}

%例题4.2.16
\begin{example}
设 $A_{3\times2},B_{2\times3}$,且
\[
AB=
\begin{pmatrix}
8&2&-2\\
2&5&4\\
-2&4&5
\end{pmatrix}.
\]
证明:
\begin{enumerate}
  \item $r(A)=r(B)=2$;
  \item $BA=9E_2$.
\end{enumerate}
\end{example}
\exampleblank

\begin{solution}
直接计算可知 $r(AB)=2$.又
\[
 r(AB)\leq r(A)\leq2,
 \qquad
 r(AB)\leq r(B)\leq2,
\]
故 $r(A)=r(B)=2$.

再计算得到
\[
 (AB)^2=9AB.
\]
于是
\[
 A(BA-9E_2)B=O.
\]
因 $A$ 列满秩,存在左逆;$B$ 行满秩,存在右逆.两边分别乘相应左逆,右逆,得
\[
 \boxed{BA=9E_2}.
\]
\end{solution}

%例题4.2.17
\begin{example}
设 $A$ 为 $n\times m$ 实矩阵,$r(A)=m$,$n\geq m$,且
\[
 (AA^{\mathsf T})^2=aAA^{\mathsf T}.
\]
证明
\[
 A^{\mathsf T}A=aE_m.
\]
\end{example}
\exampleblank

\begin{solution}
原式化为
\[
 A(A^{\mathsf T}A-aE_m)A^{\mathsf T}=O.
\]
因 $A$ 列满秩,存在左逆;$A^{\mathsf T}$ 行满秩,存在右逆.左右相乘消去 $A,A^{\mathsf T}$,即得
\[
 \boxed{A^{\mathsf T}A=aE_m}.
\]
\end{solution}

\subsection{秩一矩阵}

\begin{theorem}[\marktwostar]
设 $A$ 为 $n$ 阶矩阵且 $r(A)=1$,则:
\begin{enumerate}
  \item 存在列向量 $\alpha,\beta$,使 $A=\alpha\beta^{\mathsf T}$;
  \item $A$ 的特征值为 $0$(至少 $n-1$ 重)以及 $\operatorname{tr}A$;
  \item
  \[
  A^2=(\operatorname{tr}A)A,
  \qquad
  A^k=(\operatorname{tr}A)^{k-1}A.
  \]
\end{enumerate}
\end{theorem}

%例题4.2.18
\begin{example}
设 $A$ 为 $n$ 阶不可逆矩阵,$A^*$ 为伴随矩阵,且
\[
 \operatorname{tr}(A^*)=a\neq0.
\]
求 $r(A)$ 以及
\[
 |\lambda E-A^*|.
\]
\end{example}
\exampleblank

\begin{solution}
因 $A^*\neq O$,故 $r(A)\geq n-1$;又 $A$ 不可逆,所以
\[
 \boxed{r(A)=n-1}.
\]
此时 $r(A^*)=1$.秩一矩阵 $A^*$ 的非零特征值等于其迹 $a$,故
\[
 \boxed{|\lambda E-A^*|=\lambda^{n-1}(\lambda-a)}.
\]
\end{solution}

%例题4.2.19
\begin{example}
设 $A$ 为二阶方阵,且存在正整数 $k\geq2$ 使
\[
 A^k=O.
\]
证明 $A^2=O$.
\end{example}
\exampleblank

\begin{solution}
$A$ 为幂零矩阵,因此其全部特征值为零,故
\[
 \operatorname{tr}A=0,\qquad |A|=0.
\]
二阶矩阵的 Cayley--Hamilton 定理给出
\[
 A^2-(\operatorname{tr}A)A+|A|E=O,
\]
于是
\[
 \boxed{A^2=O}.
\]
\end{solution}


\newpage

\section{子空间的交,和与直和}

\begin{definition}
设 $V_1,V_2$ 为数域 $F$ 上的线性空间,定义
\[
V_1\oplus V_2
=
\{(\alpha,\beta)\mid \alpha\in V_1,\beta\in V_2\},
\]
并按分量定义加法和数乘,称为 $V_1,V_2$ 的外直和.
\end{definition}

若 $\alpha_1,\ldots,\alpha_m$ 与 $\beta_1,\ldots,\beta_n$ 分别为 $V_1,V_2$ 的基,则
\[
(\alpha_1,0),\ldots,(\alpha_m,0),(0,\beta_1),\ldots,(0,\beta_n)
\]
是外直和的一组基,因此
\[
\dim(V_1\oplus V_2)=\dim V_1+\dim V_2.
\]
外直和与通常的内直和在有限维情形下具有自然的同构关系.

\begin{theorem}[\markstar\ 有限不覆盖定理]
设 $V$ 是数域 $K$ 上的有限维线性空间,$V_1,\ldots,V_s$ 为有限个真子空间,则存在向量
\[
\alpha\notin V_1\cup\cdots\cup V_s.
\]
\end{theorem}

\begin{remark}
这里``数域''按通常高等代数约定为特征零的无限域.对有限域,该结论一般不成立.
\end{remark}

%例题4.3.1
\begin{example}
设 $\alpha_1,\ldots,\alpha_r$ 与 $\beta_1,\ldots,\beta_s$ 分别是 $K^n$ 中的线性无关向量组.证明:
\[
L(\alpha_1,\ldots,\alpha_r)\cap L(\beta_1,\ldots,\beta_s)
\]
的维数等于齐次方程
\[
 x_1\alpha_1+\cdots+x_r\alpha_r
 +y_1\beta_1+\cdots+y_s\beta_s=0
\]
的解空间维数.
\end{example}
\exampleblank

\begin{solution}
把一个解 $(x,y)$ 映到
\[
\displaystyle\sum_{i=1}^r x_i\alpha_i=-\sum_{j=1}^s y_j\beta_j.
\]
所得向量属于两子空间之交.由于两组向量各自线性无关,该映射既单射又满射,因此解空间与交空间线性同构,维数相同.
\end{solution}

%例题4.3.2
\begin{example}
在 $\mathbb{R}^4$ 中,令
\[
\alpha_1=(1,2,1,0),\quad
\alpha_2=(-1,1,1,1),\quad
\alpha_3=(0,3,2,1),
\]
生成子空间 $V_1$,并令
\[
\beta_1=(2,-1,0,1),\qquad
\beta_2=(1,-1,3,7)
\]
生成子空间 $V_2$.求 $V_1+V_2$ 与 $V_1\cap V_2$ 的基.
\end{example}
\exampleblank

\begin{solution}
注意
\[
\alpha_3=\alpha_1+\alpha_2,
\]
故 $V_1=L(\alpha_1,\alpha_2)$.行化简可得
\[
\{\alpha_1,\alpha_2,\beta_1\}
\]
为 $V_1+V_2$ 的一组基.

解
\[
 c_1\alpha_1+c_2\alpha_2=d_1\beta_1+d_2\beta_2
\]
得非零关系可取
\[
(c_1,c_2,d_1,d_2)=(-1,4,-3,1).
\]
故
\[
V_1\cap V_2=L((-5,2,3,4)).
\]
\end{solution}

%例题4.3.3
\begin{example}
设 $M_2(K)$ 的子空间
\[
W_1=L(A_1,A_2),\qquad W_2=L(B_1,B_2),
\]
其中
\[
A_1=\begin{pmatrix}1&1\\0&0\end{pmatrix},\quad
A_2=\begin{pmatrix}1&0\\1&1\end{pmatrix},
\]
\[
B_1=\begin{pmatrix}0&0\\1&1\end{pmatrix},\quad
B_2=\begin{pmatrix}0&1\\1&0\end{pmatrix}.
\]
求 $W_1+W_2$,$W_1\cap W_2$,并判断是否有 $M_2(K)=W_1\oplus W_2$.
\end{example}
\exampleblank

\begin{solution}
把矩阵按 $(a_{11},a_{12},a_{21},a_{22})$ 向量化.四个向量线性无关,因此
\[
\dim(W_1+W_2)=4,
\qquad W_1\cap W_2=\{O\}.
\]
故
\[
\boxed{M_2(K)=W_1\oplus W_2}.
\]
并且 $A_1,A_2,B_1,B_2$ 就是和空间的一组基.
\end{solution}

%例题4.3.4
\begin{example}
设 $W$ 为 $\mathbb{R}^n$ 的子空间,且 $W$ 中每个非零向量的零分量个数不超过 $r$.证明
\[
\dim W\leq r+1.
\]
\end{example}
\exampleblank

\begin{solution}
若 $\dim W\geq r+2$,任选 $r+1$ 个坐标位置,令这些坐标都等于零,相当于在 $W$ 上施加至多 $r+1$ 个齐次线性条件.由于 $\dim W>r+1$,必存在非零解,即存在 $W$ 中非零向量至少有 $r+1$ 个零分量,矛盾.
\end{solution}

%例题4.3.5
\begin{example}
证明:子空间 $L_1,L_2$ 的和 $S=L_1+L_2$ 是直和,当且仅当至少存在一个向量 $x\in S$ 能唯一写成
\[
x=x_1+x_2,\qquad x_i\in L_i.
\]
\end{example}
\exampleblank

\begin{solution}
若 $L_1\cap L_2=\{0\}$,则所有向量的分解都唯一.

反之,若 $L_1\cap L_2$ 中有非零向量 $u$,则任意一个分解 $x=x_1+x_2$ 都可改写为
\[
x=(x_1+u)+(x_2-u),
\]
且两项仍分别属于 $L_1,L_2$,所以没有任何向量的分解是唯一的.故只要存在一个唯一分解,就必有交为零.
\end{solution}

%例题4.3.6
\begin{example}
设 $W,W_1,W_2$ 都是 $V$ 的子空间,$W_1\subseteq W$,且
\[
V=W_1\oplus W_2.
\]
证明
\[
\dim W=\dim W_1+\dim(W_2\cap W).
\]
\end{example}
\exampleblank

\begin{solution}
任取 $w\in W$,在 $V=W_1\oplus W_2$ 中唯一写成 $w=w_1+w_2$.因为 $w,w_1\in W$,故 $w_2\in W\cap W_2$.所以
\[
W=W_1\oplus(W\cap W_2),
\]
取维数即得结论.
\end{solution}

%例题4.3.7
\begin{example}
设 $V$ 为 $n$ 维线性空间,$V_1$ 为其子空间,且
\[
\dim V_1\geq\dfrac n2.
\]
\begin{enumerate}
  \item 证明存在子空间 $W_1,W_2$,使
  \[
  V=V_1\oplus W_1=V_1\oplus W_2,
  \qquad W_1\cap W_2=\{0\};
  \]
  \item 若 $\dim V_1<n/2$,上述结论是否仍成立?
\end{enumerate}
\end{example}
\exampleblank

\begin{solution}
记 $k=\dim V_1$,$m=n-k$,则 $m\leq k$.取 $V_1$ 的基 $e_1,\ldots,e_k$,扩充为 $V$ 的基
\[
e_1,\ldots,e_k,f_1,\ldots,f_m.
\]
令
\[
W_1=L(f_1,\ldots,f_m),
\]
\[
W_2=L(f_1+e_1,\ldots,f_m+e_m).
\]
则两者都是 $V_1$ 的补空间,而且 $W_1\cap W_2=\{0\}$.

若 $k<n/2$,任意两个补空间维数均为 $n-k$,于是
\[
\dim(W_1\cap W_2)\geq2(n-k)-n=n-2k>0,
\]
故不可能交为零.
\end{solution}

%例题4.3.8
\begin{example}
设 $A,B,C,D$ 为两两可交换的 $n$ 阶矩阵,满足
\[
AC+BD=E_n.
\]
令 $W,V_1,V_2$ 分别为 $ABx=0,Bx=0,Ax=0$ 的解空间.证明
\[
W=V_1\oplus V_2.
\]
\end{example}
\exampleblank

\begin{solution}
显然 $V_1,V_2\subseteq W$.若 $x\in W$,则
\[
x=(AC+BD)x=ACx+BDx.
\]
因为各矩阵可交换且 $ABx=0$,有
\[
B(ACx)=AC(Bx)=C(ABx)=0,
\]
所以 $ACx\in V_1$;同理 $BDx\in V_2$.故 $W=V_1+V_2$.

若 $y\in V_1\cap V_2$,则 $Ay=By=0$,从而
\[
y=(AC+BD)y=0.
\]
故交为零.
\end{solution}

%例题4.3.9
\begin{example}[\markstar]
设 $M\in P^{n\times n}$,$f,g\in P[x]$ 且 $(f,g)=1$.令
\[
A=f(M),\qquad B=g(M),
\]
$W,W_1,W_2$ 分别为 $ABx=0,Ax=0,Bx=0$ 的解空间.证明
\[
W=W_1\oplus W_2.
\]
\end{example}
\exampleblank

\begin{solution}
由 Bézout 恒等式,存在 $u,v\in P[x]$ 使
\[
u(M)A+v(M)B=E.
\]
且所有这些矩阵都可交换.于是完全按上一题的证明,任意 $x\in W$ 有
\[
x=u(M)Ax+v(M)Bx,
\]
其中两项分别属于 $W_2,W_1$,且交空间为零,故结论成立.
\end{solution}

%例题4.3.10
\begin{example}
设 $m<n$,$A\in P^{m\times n}$,$B\in P^{(n-m)\times n}$.令 $V_1,V_2$ 分别为 $Ax=0,Bx=0$ 的解空间.证明
\[
P^n=V_1\oplus V_2
\]
的充分必要条件是
\[
\begin{pmatrix}A\\B\end{pmatrix}x=0
\]
只有零解.
\end{example}
\exampleblank

\begin{solution}
两方程组公共解空间正是
\[
V_1\cap V_2
=\ker\begin{pmatrix}A\\B\end{pmatrix}.
\]
若块矩阵方程只有零解,则其秩为 $n$.又
\[
r(A)\leq m,\qquad r(B)\leq n-m,
\]
所以必有 $r(A)=m,r(B)=n-m$.于是
\[
\dim V_1=n-m,\qquad \dim V_2=m,
\]
加上交为零,得到直和为 $P^n$.

反之若 $P^n=V_1\oplus V_2$,则交为零,因此块方程只有零解.
\end{solution}

%例题4.3.11
\begin{example}[\markstar]
设 $V$ 为数域 $F$ 上全体 $n$ 阶对称矩阵组成的线性空间,
\[
U=\{A\in V\mid\operatorname{tr}A=0\},
\qquad
W=\{\lambda E\mid\lambda\in F\}.
\]
\begin{enumerate}
  \item 证明 $U,W$ 为子空间;
  \item 求各自一组基和维数;
  \item 证明 $V=U\oplus W$.
\end{enumerate}
\end{example}
\exampleblank

\begin{solution}
迹为线性函数,故其核 $U$ 为子空间;$W$ 显然为一维子空间.

$V$ 的维数为 $n(n+1)/2$.$U$ 可取基
\[
E_{ij}+E_{ji}\quad(1\leq i<j\leq n),
\]
以及
\[
E_{11}-E_{nn},\ E_{22}-E_{nn},\ldots,E_{n-1,n-1}-E_{nn}.
\]
故
\[
\dim U=\dfrac{n(n+1)}2-1,
\qquad \dim W=1.
\]

对任意对称矩阵 $A$,写
\[
A=
\left(A-\dfrac{\operatorname{tr}A}{n}E\right)
+\dfrac{\operatorname{tr}A}{n}E.
\]
第一项属于 $U$,第二项属于 $W$.又 $U\cap W=\{O\}$,故
\[
\boxed{V=U\oplus W}.
\]
\end{solution}

%例题4.3.12
\begin{example}[\markstar]
设 $\mathfrak s$ 是 $M_n(F)$ 中由所有交换子
\[
AB-BA
\]
生成的子空间.证明
\[
\dim\mathfrak s=n^2-1.
\]
\end{example}
\exampleblank

\begin{solution}
任意交换子满足
\[
\operatorname{tr}(AB-BA)=0,
\]
故 $\mathfrak s$ 包含于零迹矩阵空间.

反之,零迹矩阵空间由
\[
E_{ij}\quad(i\neq j),
\qquad E_{ii}-E_{nn}\quad(i=1,\ldots,n-1)
\]
张成,而
\[
E_{ij}=[E_{ii},E_{ij}],
\qquad
E_{ii}-E_{nn}=[E_{in},E_{ni}].
\]
故两空间相等,因此
\[
\boxed{\dim\mathfrak s=n^2-1}.
\]
\end{solution}

%例题4.3.13
\begin{example}[\markstar]
设 $\mathfrak{gl}(n,F)=M_n(F)$,定义
\[
[A,B]=AB-BA.
\]
\begin{enumerate}
  \item 证明 Jacobi 恒等式
  \[
  [[A_1,A_2],A_3]+[[A_2,A_3],A_1]+[[A_3,A_1],A_2]=O;
  \]
  \item 设
  \[
  \mathfrak{sl}(n,F)=\{A\mid\operatorname{tr}A=0\},
  \]
  证明它是子空间并给出一组基;
  \item 设 $D_n$ 为数量矩阵空间,证明
  \[
  \mathfrak{gl}(n,F)=\mathfrak{sl}(n,F)\oplus D_n;
  \]
  \item 证明 $\mathfrak{sl}(n,F)$ 等于所有有限个交换子的线性和所成的空间.
\end{enumerate}
\end{example}
\exampleblank

\begin{solution}
第一问直接展开六项,逐项相消即可.

第二问中 $\mathfrak{sl}(n,F)$ 是迹映射的核,因此为子空间.一组基为
\[
E_{ij}\ (i\neq j),
\qquad
E_{ii}-E_{nn}\ (i=1,\ldots,n-1).
\]

第三问在特征零数域上,任意 $A$ 有唯一分解
\[
A=
\left(A-\dfrac{\operatorname{tr}A}{n}E\right)
+\dfrac{\operatorname{tr}A}{n}E,
\]
故为直和.

第四问由上一题中基向量均可写成交换子立即得到.
\end{solution}

%例题4.3.14
\begin{example}
设 $W_1,\ldots,W_s$ 是 $V$ 的子空间.证明
\[
\displaystyle\sum\limits_{i=1}^sW_i
\]
为直和的充要条件是
\[
W_i\cap\sum\limits_{k=1}^{i-1}W_k=\{0\},
\qquad i=2,\ldots,s.
\]
\end{example}
\exampleblank

\begin{solution}
若为直和,任一 $W_i$ 与其余各项之和的交都只能为零,故必要性成立.

反之,若逐项满足上述交为零,则对
\[
w_1+\cdots+w_s=0,
\qquad w_i\in W_i,
\]
从最后一项开始,有
\[
w_s\in W_s\cap(W_1+\cdots+W_{s-1})=\{0\}.
\]
依次倒推得到所有 $w_i=0$,故分解唯一,即为直和.
\end{solution}

%例题4.3.15
\begin{example}
设 $U_1,U_2$ 是 $n$ 维空间 $V$ 的子空间,且
\[
\dim U_1\leq m,\qquad \dim U_2\leq m<n.
\]
证明存在子空间 $W$,满足
\[
\dim W=n-m,
\qquad W\cap U_1=W\cap U_2=\{0\}.
\]
\end{example}
\exampleblank

\begin{solution}
利用有限不覆盖定理归纳选取 $w_1,\ldots,w_{n-m}$.已选 $k<n-m$ 个时,两个空间
\[
U_i+L(w_1,\ldots,w_k)
\]
的维数至多为 $m+k<n$,均为真子空间.可选 $w_{k+1}$ 同时不落在这两个子空间中.

最后令 $W=L(w_1,\ldots,w_{n-m})$.构造保证 $w_i$ 线性无关,且 $W$ 与两个 $U_i$ 的交都为零.
\end{solution}

%例题4.3.16
\begin{example}
证明:$V$ 的两个子空间 $V_1,V_2$ 的并集 $V_1\cup V_2$ 是子空间,当且仅当
\[
V_1\subseteq V_2
\quad\text{或}\quad
V_2\subseteq V_1.
\]
\end{example}
\exampleblank

\begin{solution}
充分性显然.反之若二者互不包含,取
\[
x\in V_1\setminus V_2,
\qquad
y\in V_2\setminus V_1.
\]
若并集是子空间,则 $x+y\in V_1\cup V_2$.若 $x+y\in V_1$,则 $y=(x+y)-x\in V_1$,矛盾;另一种情形同理.
\end{solution}

%例题4.3.17
\begin{example}
设 $V$ 为数域 $K$ 上的 $n$ 维线性空间.证明:对任意自然数 $m>n$,存在 $V$ 中 $m$ 个向量,使其中任意 $n$ 个都构成 $V$ 的一组基.
\end{example}
\exampleblank

\begin{solution}
取 $V$ 的一组基,把 $V$ 与 $K^n$ 同一.因数域 $K$ 为无限域,可取两两不同的 $t_1,\ldots,t_m\in K$,令
\[
v(t_j)=(1,t_j,t_j^2,\ldots,t_j^{n-1})^{\mathsf T}.
\]
任取其中 $n$ 个向量所成矩阵都是 Vandermonde 矩阵,其行列式为
\[
\displaystyle\prod_{i<j}(t_j-t_i)\neq0.
\]
故任意 $n$ 个都为一组基.
\end{solution}


\newpage

\section{线性方程组解的结构}

\subsection{齐次线性方程组}

\begin{definition}
齐次方程组 $Ax=0$ 的全部解构成一个线性子空间,称为解空间.非零解空间的一组基称为该方程组的基础解系.
\end{definition}

\begin{theorem}[\marktwostar]
设 $A$ 为 $m\times n$ 矩阵,$r(A)=r$.则 $Ax=0$ 有非零解的充要条件是 $r<n$.此时解空间维数为 $n-r$,若 $\xi_1,\ldots,\xi_{n-r}$ 为一组基础解系,则通解为
\[
x=k_1\xi_1+\cdots+k_{n-r}\xi_{n-r}.
\]
\end{theorem}

%例题4.4.1
\begin{example}
求 $\lambda$,使齐次方程组
\[
\begin{cases}
(\lambda+3)x_1+x_2+2x_3=0,\\
\lambda x_1+(\lambda-1)x_2+x_3=0,\\
3(\lambda+1)x_1+\lambda x_2+(\lambda+3)x_3=0
\end{cases}
\]
有非零解,并求一般解.
\end{example}
\exampleblank

\begin{solution}
系数行列式为
\[
\lambda^2(\lambda-1).
\]
故只需考虑 $\lambda=0,1$.

当 $\lambda=0$ 时,基础解系可取
\[
(-1,1,1)^{\mathsf T},
\]
所以
\[
x=t(-1,1,1)^{\mathsf T}.
\]
当 $\lambda=1$ 时,基础解系可取
\[
(-1,2,1)^{\mathsf T},
\]
所以
\[
x=t(-1,2,1)^{\mathsf T}.
\]
\end{solution}

%例题4.4.2
\begin{example}
设齐次方程组
\[
\begin{cases}
x_2+ax_3+bx_4=0,\\
-x_1+cx_3+dx_4=0,\\
ax_1+cx_2-ex_4=0,\\
bx_1+dx_2-ex_3=0
\end{cases}
\]
的一般解以 $x_3,x_4$ 为自由未知量.
\begin{enumerate}
  \item 求 $a,b,c,d,e$ 满足的条件;
  \item 求基础解系.
\end{enumerate}
\end{example}
\exampleblank

\begin{solution}
由前两式
\[
x_1=cx_3+dx_4,
\qquad
x_2=-ax_3-bx_4.
\]
代入后两式,要求对任意 $x_3,x_4$ 恒成立,得到
\[
ad-bc-e=0,
\qquad
bc-ad-e=0.
\]
故
\[
\boxed{e=0,
\qquad ad=bc}.
\]
此时基础解系可取
\[
\boxed{
(c,-a,1,0)^{\mathsf T},
\quad
(d,-b,0,1)^{\mathsf T}
}.
\]
\end{solution}

%例题4.4.3
\begin{example}
设四元齐次方程组 $(I)$ 为
\[
\begin{cases}
x_1+x_3=0,\\
x_2-x_4=0,
\end{cases}
\]
而方程组 $(II)$ 的通解为
\[
k_1(0,1,1,0)^{\mathsf T}
+k_2(-1,2,2,1)^{\mathsf T}.
\]
\begin{enumerate}
  \item 求 $(I)$ 的基础解系;
  \item 求两方程组的全部非零公共解.
\end{enumerate}
\end{example}
\exampleblank

\begin{solution}
$(I)$ 的通解为
\[
(-x_3,x_4,x_3,x_4)^{\mathsf T},
\]
故基础解系可取
\[
(-1,0,1,0)^{\mathsf T},
\qquad
(0,1,0,1)^{\mathsf T}.
\]
求两个解空间的交,得到一维空间
\[
L((-1,1,1,1)^{\mathsf T}).
\]
因此所有非零公共解为
\[
\boxed{t(-1,1,1,1)^{\mathsf T},\qquad t\neq0}.
\]
\end{solution}

%例题4.4.4
\begin{example}
设四元齐次方程组 $(I)$ 为
\[
\begin{cases}
2x_1+3x_2-x_3=0,\\
x_1+2x_2+x_3-x_4=0,
\end{cases}
\]
另一方程组 $(II)$ 的基础解系为
\[
\alpha_1=(2,-1,a+2,1)^{\mathsf T},
\qquad
\alpha_2=(-1,2,4,a+8)^{\mathsf T}.
\]
\begin{enumerate}
  \item 求 $(I)$ 的基础解系;
  \item 求 $a$,使 $(I),(II)$ 有非零公共解,并求全部公共解.
\end{enumerate}
\end{example}
\exampleblank

\begin{solution}
$(I)$ 的基础解系可取
\[
\eta_1=(5,-3,1,0)^{\mathsf T},
\qquad
\eta_2=(-3,2,0,1)^{\mathsf T}.
\]
四个向量 $\eta_1,\eta_2,\alpha_1,\alpha_2$ 的行列式为
\[
(a+1)^2.
\]
故有非零公共解当且仅当
\[
\boxed{a=-1}.
\]
此时直接代入可知 $\alpha_1,\alpha_2$ 都满足 $(I)$,而两者线性无关,因此两个解空间实际上相同.全部公共解为
\[
\boxed{s(2,-1,1,1)^{\mathsf T}+t(-1,2,4,7)^{\mathsf T}}.
\]
\end{solution}

%例题4.4.5
\begin{example}
已知齐次方程组
\[
(I):\quad
\begin{cases}
x_1+2x_2+3x_3=0,\\
2x_1+3x_2+5x_3=0,\\
x_1+x_2+ax_3=0,
\end{cases}
\]
和
\[
(II):\quad
\begin{cases}
x_1+bx_2+cx_3=0,\\
2x_1+b^2x_2+(c+1)x_3=0
\end{cases}
\]
同解,求 $a,b,c$.
\end{example}
\exampleblank

\begin{solution}
$(II)$ 是两方程三未知量,必有非零解,所以 $(I)$ 也必须有非零解.由
\[
\det
\begin{pmatrix}
1&2&3\\2&3&5\\1&1&a
\end{pmatrix}
=2-a
\]
得 $a=2$.此时 $(I)$ 的解空间为
\[
L((-1,-1,1)^{\mathsf T}).
\]
要使 $(II)$ 同解,两行都必须与该向量正交:
\[
-1-b+c=0,
\qquad
-2-b^2+c+1=0.
\]
故 $c=1+b=1+b^2$,即 $b=0$ 或 $1$.$b=0$ 时两行相关,解空间维数为 $2$,舍去.因此
\[
\boxed{a=2,\quad b=1,\quad c=2}.
\]
\end{solution}

%例题4.4.6
\begin{example}
已知齐次线性方程组
\[
\begin{cases}
(a_1+b)x_1+a_2x_2+\cdots+a_nx_n=0,\\
a_1x_1+(a_2+b)x_2+\cdots+a_nx_n=0,\\
\qquad\vdots\\
a_1x_1+a_2x_2+\cdots+(a_n+b)x_n=0,
\end{cases}
\]
其中 $\displaystyle\prod_{i=1}^na_i\neq0$.讨论何时只有零解,何时有非零解,并给出基础解系.
\end{example}
\exampleblank

\begin{solution}
系数矩阵为
\[
A=bE+\boldsymbol{1}a^{\mathsf T},
\qquad
a=(a_1,\ldots,a_n)^{\mathsf T}.
\]
其行列式为
\[
|A|=b^{n-1}\left(b+\sum\limits_{i=1}^na_i\right).
\]
因此若
\[
b\neq0,
\qquad
b+\sum a_i\neq0,
\]
则只有零解.

若 $b=0$,方程组退化为
\[
\displaystyle\sum a_ix_i=0,
\]
解空间维数 $n-1$,可取基础解系
\[
a_{j+1}e_j-a_je_{j+1},
\qquad j=1,\ldots,n-1.
\]

若 $b\neq0$ 且
\[
b+\sum a_i=0,
\]
则
\[
A\boldsymbol{1}=0,
\]
且秩为 $n-1$,故基础解系为
\[
\boxed{(1,1,\ldots,1)^{\mathsf T}}.
\]
\end{solution}

%例题4.4.7
\begin{example}
设 $A,B$ 为 $n$ 阶方阵,且
\[
r(A)+r(B)<n.
\]
证明 $Ax=0$ 与 $Bx=0$ 有非零公共解.
\end{example}
\exampleblank

\begin{solution}
设 $V_A=\ker A,V_B=\ker B$.则
\[
\dim V_A+\dim V_B
=2n-r(A)-r(B)>n.
\]
由维数公式,
\[
\dim(V_A\cap V_B)>0,
\]
故存在非零公共解.
\end{solution}

%例题4.4.8
\begin{example}
设 $AB=BA$,证明
\[
r(A+B)\leq r(A)+r(B)-r(AB).
\]
\end{example}
\exampleblank

\begin{solution}
显然
\[
\operatorname{Im}(A+B)\subseteq\operatorname{Im}A+\operatorname{Im}B.
\]
又因 $AB=BA$,
\[
\operatorname{Im}(AB)\subseteq\operatorname{Im}A\cap\operatorname{Im}B.
\]
因此
\[
\begin{aligned}
r(A+B)
&\leq\dim(\operatorname{Im}A+\operatorname{Im}B)\\
&=r(A)+r(B)-\dim(\operatorname{Im}A\cap\operatorname{Im}B)\\
&\leq r(A)+r(B)-r(AB).
\end{aligned}
\]
\end{solution}

%例题4.4.9
\begin{example}
设 $\eta_1,\ldots,\eta_{n-r}$ 是 $Ax=0$ 的基础解系,
\[
B=(\eta_1,\ldots,\eta_{n-r}).
\]
若 $AC=O$,证明存在唯一矩阵 $D$,使
\[
C=BD.
\]
\end{example}
\exampleblank

\begin{solution}
$AC=O$ 表明 $C$ 的每一列都属于 $\ker A$.由于 $B$ 的列向量是 $\ker A$ 的一组基,每一列都能唯一表示成 $B$ 的列向量线性组合.把这些唯一坐标列排在一起即得唯一矩阵 $D$,满足 $C=BD$.
\end{solution}

\begin{proposition}[\markstar]
设 $A,B$ 都是 $m\times n$ 矩阵,$W_A,W_B$ 分别为 $Ax=0,Bx=0$ 的解空间.则:
\begin{enumerate}
  \item $W_A\subseteq W_B$ 当且仅当存在 $m$ 阶矩阵 $P$ 使 $PA=B$;
  \item $W_A=W_B$ 当且仅当 $A,B$ 的行空间相同;在两矩阵行秩相同并取相同行数表示时,可由可逆行变换互相得到.
\end{enumerate}
\end{proposition}

%例题4.4.10
\begin{example}
设 $V_0=\ker(AB)$,求
\[
BV_0=\{Bx\mid x\in V_0\}
\]
的维数.
\end{example}
\exampleblank

\begin{solution}
考虑 $B$ 在 $V_0$ 上的限制.其核为
\[
V_0\cap\ker B=\ker B,
\]
因为 $\ker B\subseteq\ker(AB)$.故
\[
\dim BV_0
=\dim V_0-\dim\ker B
=(n-r(AB))-(n-r(B)).
\]
所以
\[
\boxed{\dim BV_0=r(B)-r(AB)}.
\]
\end{solution}

%例题4.4.11
\begin{example}
设
\[
A=\begin{pmatrix}A_1\\A_2\end{pmatrix}
\]
为数域 $P$ 上的 $n$ 阶方阵,其中 $A_1,A_2$ 分别为 $k\times n$ 与 $(n-k)\times n$ 矩阵.令 $V_1,V_2$ 分别为 $A_1x=0,A_2x=0$ 的解空间.原题要求证明
\[
P^n=V_1\oplus V_2
\iff
r(A)=r(A_1)+r(A_2).
\]
\end{example}
\exampleblank

\begin{solution}
原命题按字面并不成立.例如取 $A_1=A_2=O$,则右端成立,但 $V_1=V_2=P^n$,不可能为直和.

由维数公式可得
\[
\dim(V_1+V_2)
=n+r(A)-r(A_1)-r(A_2).
\]
因此原右端实际上等价于
\[
V_1+V_2=P^n,
\]
而不是直和.

若要求
\[
P^n=V_1\oplus V_2,
\]
则必须且只需
\[
\ker A=V_1\cap V_2=\{0\},
\]
也就是
\[
\boxed{r(A)=n}.
\]
此时因 $A$ 的 $n$ 行线性无关,自然有 $r(A_1)=k,r(A_2)=n-k$,从而右端秩等式也成立.
\end{solution}

%例题4.4.12
\begin{example}
设 $A$ 为 $(n-1)\times n$ 矩阵,划去第 $j$ 列所得的 $(n-1)$ 阶子式为 $\mu_j$.证明:
\begin{enumerate}
  \item
  \[
  (\mu_1,-\mu_2,\ldots,(-1)^{n-1}\mu_n)^{\mathsf T}
  \]
  是 $Ax=0$ 的一个解;
  \item 若 $r(A)=n-1$,则所有解都是该向量的倍数.
\end{enumerate}
\end{example}
\exampleblank

\begin{solution}
设
\[
\eta=(\mu_1,-\mu_2,\ldots,(-1)^{n-1}\mu_n)^{\mathsf T}.
\]
对 $A$ 的任一行与 $\eta$ 作内积,正好等于把该行重复添加成第 $n$ 行所得 $n$ 阶行列式按最后一行展开,因此为零.故 $A\eta=0$.

若 $r(A)=n-1$,则至少有一个 $\mu_j\neq0$,所以 $\eta\neq0$;同时 $\ker A$ 为一维,故所有解都是 $\eta$ 的倍数.
\end{solution}

%例题4.4.13
\begin{example}
设 $A,B$ 都是数域 $F$ 上的 $(n-1)\times n$ 行满秩矩阵.记 $a_i,b_i$ 为分别划去第 $i$ 列所得的 $(n-1)$ 阶子式.若 $Ax=0$ 与 $Bx=0$ 同解,证明向量
\[
(a_1,a_2,\ldots,a_n),
\qquad
(b_1,b_2,\ldots,b_n)
\]
的对应分量成比例.
\end{example}
\exampleblank

\begin{solution}
由上一题,$\ker A$ 与 $\ker B$ 分别由
\[
(a_1,-a_2,\ldots,(-1)^{n-1}a_n)^{\mathsf T}
\]
和
\[
(b_1,-b_2,\ldots,(-1)^{n-1}b_n)^{\mathsf T}
\]
张成.两核相同且均为一维,因此两个向量成比例.消去相同的交替符号,即得 $(a_i)$ 与 $(b_i)$ 对应分量成比例.
\end{solution}

%例题4.4.14
\begin{example}
设 $(n-1)\times n$ 矩阵 $A$ 的每个行向量都满足
\[
x_1+x_2+\cdots+x_n=0.
\]
记划去第 $i$ 列所得子式为 $\mu_i$.
\begin{enumerate}
  \item 证明
  \[
  \displaystyle\sum\limits_{i=1}^n(-1)^i\mu_i=0
  \]
  当且仅当 $A$ 的行向量不是该超平面的基础解系;
  \item 若
  \[
  \displaystyle\sum\limits_{i=1}^n(-1)^i\mu_i=1,
  \]
  求 $\mu_i$.
\end{enumerate}
\end{example}
\exampleblank

\begin{solution}
因为每一行坐标和为零,故
\[
A\boldsymbol{1}=0.
\]
若 $r(A)=n-1$,则上一题中的余子式向量非零且张成 $\ker A$,因此存在 $c\neq0$ 使
\[
(\mu_1,-\mu_2,\ldots,(-1)^{n-1}\mu_n)=c(1,1,\ldots,1).
\]
即
\[
\mu_i=c(-1)^{i-1}.
\]
于是
\[
\displaystyle\sum_{i=1}^n(-1)^i\mu_i=-nc\neq0.
\]
反之若 $r(A)<n-1$,所有 $\mu_i=0$,该和为零.因此第一问成立.

第二问有
\[
-nc=1,
\]
故
\[
\boxed{\mu_i=\dfrac{(-1)^i}{n}}.
\]
\end{solution}

\subsection{非齐次线性方程组}

\begin{theorem}[\marktwostar]
对非齐次方程组 $Ax=\beta$,下列命题等价:
\begin{enumerate}
  \item 方程有解;
  \item $\beta$ 属于 $A$ 的列空间;
  \item $r(A)=r(A,\beta)$.
\end{enumerate}
若 $\eta_0$ 是一个特解,而 $\xi_1,\ldots,\xi_{n-r}$ 是导出组 $Ax=0$ 的基础解系,则通解为
\[
x=\eta_0+k_1\xi_1+\cdots+k_{n-r}\xi_{n-r}.
\]
\end{theorem}

%例题4.4.15
\begin{example}
求解下列方程组:
\begin{enumerate}
\item
\[
\begin{cases}
2x_1+x_2+x_3=2,\\
x_1+3x_2+x_3=5,\\
x_1+x_2+5x_3=-7,\\
2x_1+3x_2-3x_3=14;
\end{cases}
\]
\item
\[
\begin{cases}
2x_1+7x_2+3x_3+x_4=5,\\
x_1+3x_2+5x_3-2x_4=3,\\
x_1+5x_2-9x_3+8x_4=1,\\
5x_1+18x_2+4x_3+5x_4=12;
\end{cases}
\]
\item
\[
\begin{cases}
2x_1-x_2+x_3-x_4=3,\\
4x_1-2x_2-2x_3+3x_4=2,\\
2x_1-x_2+5x_3-6x_4=1,\\
2x_1-x_2-3x_3+4x_4=5;
\end{cases}
\]
\item
\[
\begin{cases}
x_1+x_2+x_3+x_4+x_5=1,\\
3x_1+2x_2+x_3+x_4-3x_5=-2,\\
x_2+2x_3+2x_4+6x_5=5,\\
5x_1+4x_2+3x_3+3x_4-x_5=0.
\end{cases}
\]
\end{enumerate}
\end{example}
\exampleblank

\begin{solution}
逐一行化简得:
\begin{enumerate}
  \item 唯一解
  \[
  \boxed{(x_1,x_2,x_3)=(1,2,-2)};
  \]
  \item 令 $x_3=s,x_4=t$,则
  \[
  \boxed{
  x_1=-26s+17t+6,
  \quad x_2=7s-5t-1
  };
  \]
  \item 系数矩阵秩为 $2$,增广矩阵秩为 $3$,故
  \[
  \boxed{\text{无解}};
  \]
  \item 令 $x_3=s,x_4=t,x_5=u$,则
  \[
  \boxed{
  x_1=s+t+5u-4,
  \quad
  x_2=-2s-2t-6u+5
  }.
  \]
\end{enumerate}
\end{solution}

%例题4.4.16
\begin{example}
设 $\alpha_1,\alpha_2,\alpha_3$ 是方程组
\[
\begin{cases}
x-3y+z=2,\\2x+y+tz=-1,\\7x-2z=-1\end{cases}
\]
的三个互异解向量.
\begin{enumerate}
  \item 求 $t$;
  \item 证明 $\alpha_1-\alpha_2$ 与 $\alpha_1-\alpha_3$ 线性相关.
\end{enumerate}
\end{example}
\exampleblank

\begin{solution}
存在三个互异解说明解集不是单点,因此系数矩阵奇异.其行列式为
\[
-21(t+1),
\]
故
\[
\boxed{t=-1}.
\]
此时系数矩阵秩为 $2$,导出齐次方程组解空间维数为 $1$.任意两个特解之差都属于这一维空间,因此 $\alpha_1-\alpha_2$ 与 $\alpha_1-\alpha_3$ 必线性相关.
\end{solution}

%例题4.4.17
\begin{example}
讨论参数 $a,b$,使方程组
\[
\begin{cases}
ax_1+(b+1)x_2+2x_3=1,\\
ax_1+(2b+1)x_2+3x_3=1,\\
ax_1+(b+1)x_2+(b+4)x_3=2b+1
\end{cases}
\]
有解,并求其解.
\end{example}
\exampleblank

\begin{solution}
第二式减第一式,第三式减第一式,得到
\[
bx_2+x_3=0,
\qquad
(b+2)x_3=2b.
\]
分情况:

若 $b=-2$,第二个式子化为 $0=-4$,故无解.

若 $b=0$,则 $x_3=0$,原方程化为
\[
ax_1+x_2=1,
\]
故对任意 $a$ 均有无穷多解:
\[
\boxed{x_1=t,\quad x_2=1-at,\quad x_3=0}.
\]

若 $b\neq0,-2$,则
\[
x_3=\dfrac{2b}{b+2},
\qquad
x_2=-\dfrac{2}{b+2}.
\]
第一式给出
\[
a x_1=\dfrac{4-b}{b+2}.
\]
若 $a\neq0$,有唯一解
\[
\boxed{
 x_1=\dfrac{4-b}{a(b+2)},
 \quad x_2=-\dfrac{2}{b+2},
 \quad x_3=\dfrac{2b}{b+2}
 }.
\]
若 $a=0$,仅当 $b=4$ 时有解,此时
\[
\boxed{x_1=t,\quad x_2=-\dfrac13,\quad x_3=\dfrac43}.
\]
其余 $a=0$ 的情形无解.
\end{solution}

%例题4.4.18
\begin{example}
设 $A$ 为三对角矩阵
\[
A=
\begin{pmatrix}
2a&1&0&\cdots&0\\
a^2&2a&1&\ddots&\vdots\\
0&a^2&2a&\ddots&0\\
\vdots&\ddots&\ddots&\ddots&1\\
0&\cdots&0&a^2&2a
\end{pmatrix}_{n\times n},
\]
且 $Ax=e_1$.
\begin{enumerate}
  \item 证明 $|A|=(n+1)a^n$;
  \item 何时有唯一解,并求 $x_1$;
  \item 何时有无穷多解,并求通解.
\end{enumerate}
\end{example}
\exampleblank

\begin{solution}
记 $D_n=|A|$,则
\[
D_n=2aD_{n-1}-a^2D_{n-2},
\qquad
D_0=1,
\quad D_1=2a.
\]
归纳得
\[
\boxed{D_n=(n+1)a^n}.
\]

因此 $a\neq0$ 时有唯一解.由 Cramer 法则,
\[
x_1=\dfrac{D_{n-1}}{D_n}
=\boxed{\dfrac{n}{(n+1)a}}.
\]

当 $a=0$ 时,矩阵只在超对角线上有 $1$,方程组为
\[
x_2=1,\quad x_3=\cdots=x_n=0,
\]
而 $x_1$ 自由,因此
\[
\boxed{x=(t,1,0,\ldots,0)^{\mathsf T}}.
\]
\end{solution}

%例题4.4.19
\begin{example}
讨论参数 $a,b$,使方程组
\[
\begin{cases}
x_1+x_2+x_3+x_4+x_5=1,\\
3x_1+2x_2+x_3+x_4-3x_5=a,\\
x_2+2x_3+2x_4+6x_5=3,\\
5x_1+4x_2+3x_3+3x_4-x_5=b
\end{cases}
\]
有解,并求一般解.
\end{example}
\exampleblank

\begin{solution}
系数矩阵的两个独立行关系为
\[
-3R_1+R_2+R_3=0,
\qquad
-2R_1-R_2+R_4=0.
\]
故常数项必须满足
\[
-3+a+3=0,
\qquad
-2-a+b=0.
\]
于是
\[
\boxed{a=0,\qquad b=2}.
\]
此时令 $x_3=s,x_4=t,x_5=u$,通解为
\[
\boxed{
 x_1=s+t+5u-2,
 \quad
 x_2=-2s-2t-6u+3,
 \quad
 x_3=s,
 \quad x_4=t,
 \quad x_5=u
 }.
\]
\end{solution}

%例题4.4.20
\begin{example}
已知
\[
\alpha_1=(1,0,2,3)^{\mathsf T},
\quad
\alpha_2=(1,1,3,5)^{\mathsf T},
\]
\[
\alpha_3=(1,-1,a+2,1)^{\mathsf T},
\quad
\alpha_4=(1,2,4,a+8)^{\mathsf T},
\]
\[
\beta=(1,1,b+3,5)^{\mathsf T}.
\]
讨论何时 $\beta$ 能或不能由 $\alpha_1,\ldots,\alpha_4$ 线性表示,并写出表示式.
\end{example}
\exampleblank

\begin{solution}
有
\[
\det(\alpha_1,\alpha_2,\alpha_3,\alpha_4)=(a+1)^2.
\]
若 $a\neq-1$,四向量为一组基,因此任意 $b$ 都可表示,且解为
\[
\boxed{
\beta
=-\dfrac{2b}{a+1}\alpha_1
+\dfrac{a+b+1}{a+1}\alpha_2
+\dfrac{b}{a+1}\alpha_3
}.
\]

若 $a=-1$,四向量张成二维子空间.代入检验可得 $\beta$ 属于该子空间当且仅当 $b=0$.因此
\[
\boxed{a=-1,\ b\neq0}
\]
时不能表示;而 $a=-1,b=0$ 时有无穷多种表示,例如
\[
\boxed{\beta=\alpha_2}.
\]
\end{solution}

%例题4.4.21
\begin{example}
设 $A=(\alpha_1,\ldots,\alpha_n)$ 的前 $n-1$ 个列向量线性相关,而后 $n-1$ 个列向量线性无关,令
\[
\beta=\alpha_1+\cdots+\alpha_n.
\]
\begin{enumerate}
  \item 证明 $Ax=\beta$ 有无穷多解;
  \item 求通解的形式.
\end{enumerate}
\end{example}
\exampleblank

\begin{solution}
后 $n-1$ 列线性无关,故 $r(A)\geq n-1$;前 $n-1$ 列线性相关说明 $A$ 不可能满秩,因此
\[
r(A)=n-1.
\]
而
\[
A\boldsymbol{1}=\beta,
\]
故方程有解且导出组有一维非零解空间,所以有无穷多解.

因 $\alpha_2,\ldots,\alpha_n$ 线性无关,可唯一写
\[
\alpha_1=c_2\alpha_2+\cdots+c_n\alpha_n.
\]
于是
\[
\eta=(1,-c_2,\ldots,-c_n)^{\mathsf T}
\]
张成 $\ker A$.通解为
\[
\boxed{x=\boldsymbol{1}+t\eta}.
\]
\end{solution}

%例题4.4.22
\begin{example}
设 $A=(A_1,\ldots,A_n)$,方程 $Ax=b$ 的通解为
\[
x=\eta_0+k_1\xi_1+\cdots+k_s\xi_s,
\]
其中
\[
\eta_0=(1,1,\ldots,1)^{\mathsf T},
\]
而 $\xi_i$ 的前 $i$ 个分量为 $1$,其余为 $0$.又设
\[
B=(nA_n,(n-1)A_{n-1},\ldots,2A_2,A_1).
\]
求 $Bx=b$ 的通解.
\end{example}
\exampleblank

\begin{solution}
定义可逆线性变换 $T$ 使
\[
(Tx)_j=jx_{n-j+1}.
\]
则
\[
Bx=A(Tx).
\]
所以 $Bx=b$ 当且仅当 $y=Tx$ 是 $Ay=b$ 的解.

由 $x=T^{-1}y$,特解变为
\[
\eta_0'=\left(\dfrac1n,\dfrac1{n-1},\ldots,\dfrac12,1\right)^{\mathsf T}.
\]
而 $T^{-1}\xi_i$ 是最后 $i$ 个分量分别为
\[
\dfrac1i,\dfrac1{i-1},\ldots,1
\]
,其余为零的向量,记为 $\zeta_i$.故
\[
\boxed{x=\eta_0'+k_1\zeta_1+\cdots+k_s\zeta_s}.
\]
\end{solution}

%例题4.4.23
\begin{example}
设 $A$ 为 $n$ 阶方阵,$A^*=(A_{ij})$ 为伴随矩阵且 $A_{11}\neq0$,又设 $b\neq0$.证明:
\[
Ax=b\text{ 有无穷多解}
\iff
A^*b=0.
\]
\end{example}
\exampleblank

\begin{solution}
因为 $A_{11}\neq0$,只要 $|A|=0$ 就有 $r(A)=n-1$,从而 $r(A^*)=1$.

若 $Ax=b$ 有无穷多解,则 $|A|=0$ 且 $b\in\operatorname{Im}A$.由
\[
A^*A=|A|E=O
\]
得 $\operatorname{Im}A\subseteq\ker A^*$,故 $A^*b=0$.

反之若 $A^*b=0$.若 $|A|\neq0$,则 $A^*$ 可逆,会推出 $b=0$,矛盾.因此 $|A|=0$,且由 $A_{11}\neq0$ 得 $r(A)=n-1$.此时
\[
\dim\ker A^*=n-1=\dim\operatorname{Im}A,
\]
结合 $\operatorname{Im}A\subseteq\ker A^*$ 得二者相等,所以 $b\in\operatorname{Im}A$,方程有解;又 $\ker A$ 一维,故有无穷多解.
\end{solution}

%例题4.4.24
\begin{example}
设 $A$ 为 $n$ 阶方阵,$A^*$ 为伴随矩阵,且存在非零列向量 $\alpha$ 满足
\[
A\alpha=0.
\]
证明
\[
A^*x=\alpha\text{ 有解}
\iff
r(A)=n-1.
\]
\end{example}
\exampleblank

\begin{solution}
若 $r(A)\leq n-2$,则 $A^*=O$,而 $\alpha\neq0$,故方程无解.

若 $r(A)=n-1$,则 $r(A^*)=1$.又
\[
AA^*=O,
\]
故
\[
\operatorname{Im}A^*\subseteq\ker A.
\]
两边维数都为 $1$,所以二者相等.因 $\alpha\in\ker A$ 且非零,必有 $\alpha\in\operatorname{Im}A^*$,故方程有解.
\end{solution}

\subsection{对角占优矩阵}

\begin{proposition}[\markstar]
设 $A=(a_{ij})\in M_n(\mathbb{R})$.
\begin{enumerate}
  \item 若
  \[
  |a_{ii}|>\sum\limits_{j\neq i}|a_{ij}|,
  \qquad i=1,\ldots,n,
  \]
  则 $|A|\neq0$;
  \item 若
  \[
  a_{ii}>\sum\limits_{j\neq i}|a_{ij}|,
  \qquad i=1,\ldots,n,
  \]
  则 $|A|>0$.
\end{enumerate}
\end{proposition}

\begin{proof}
第一问若 $Ax=0$ 有非零解,取 $|x_k|=\max_i|x_i|>0$,则
\[
|a_{kk}||x_k|
=\left|\sum_{j\neq k}a_{kj}x_j\right|
\leq\sum_{j\neq k}|a_{kj}||x_j|
\leq\left(\sum_{j\neq k}|a_{kj}|\right)|x_k|,
\]
与严格对角占优矛盾,故 $A$ 可逆.

第二问考虑连续矩阵族
\[
A(t)=D+t(A-D),\qquad 0\leq t\leq1,
\]
其中 $D=\diag(a_{11},\ldots,a_{nn})$.每个 $A(t)$ 仍严格对角占优,所以行列式从不为零.又
\[
|A(0)|=\prod_i a_{ii}>0,
\]
故连续性给出 $|A(1)|=|A|>0$.
\end{proof}

\subsection{行(列)和相等}

%例题4.4.25
\begin{example}
设 $A$ 为可逆矩阵,且每行元素之和均为 $a$.证明对任意正整数 $k$,$A^{-k}$ 的每行元素之和均为 $a^{-k}$.
\end{example}
\exampleblank

\begin{solution}
令 $\boldsymbol{1}=(1,\ldots,1)^{\mathsf T}$,则
\[
A\boldsymbol{1}=a\boldsymbol{1}.
\]
因 $A$ 可逆,必有 $a\neq0$,且
\[
A^{-1}\boldsymbol{1}=a^{-1}\boldsymbol{1}.
\]
迭代得到
\[
A^{-k}\boldsymbol{1}=a^{-k}\boldsymbol{1}.
\]
这正表示每行元素之和为 $a^{-k}$.
\end{solution}

%例题4.4.26
\begin{example}
设 $A$ 为行等和矩阵.证明:
\begin{enumerate}
  \item 对任意正整数 $k$,$A^k$ 仍为行等和矩阵;
  \item 若 $A$ 可逆,则 $A^{-1}$ 仍为行等和矩阵.
\end{enumerate}
\end{example}
\exampleblank

\begin{solution}
若 $A\boldsymbol{1}=a\boldsymbol{1}$,则
\[
A^k\boldsymbol{1}=a^k\boldsymbol{1}.
\]
若 $A$ 可逆,则 $a\neq0$,且
\[
A^{-1}\boldsymbol{1}=a^{-1}\boldsymbol{1}.
\]
故结论成立.
\end{solution}

%例题4.4.27
\begin{example}
证明:若可逆矩阵 $A$ 的各行元素之和相等,则 $A$ 中各列元素的代数余子式之和也相等.
\end{example}
\exampleblank

\begin{solution}
设 $A\boldsymbol{1}=a\boldsymbol{1}$.由可逆性 $a\neq0$,所以
\[
A^{-1}\boldsymbol{1}=a^{-1}\boldsymbol{1}.
\]
又
\[
A^*=|A|A^{-1},
\]
故
\[
A^*\boldsymbol{1}=|A|a^{-1}\boldsymbol{1}.
\]
$A^*$ 的第 $i$ 行正由 $A$ 的第 $i$ 列代数余子式组成,因此上式恰说明各列代数余子式之和相等.
\end{solution}

%例题4.4.28
\begin{example}
设 $A=(a_{ij})_{n\times n}$ 且每行元素之和为零,即
\[
\displaystyle\sum\limits_{j=1}^na_{ij}=0,
\qquad i=1,\ldots,n.
\]
证明第一行各元素的代数余子式相等:
\[
A_{11}=A_{12}=\cdots=A_{1n}.
\]
\end{example}
\exampleblank

\begin{solution}
条件等价于
\[
A\boldsymbol{1}=0.
\]
由
\[
A^*A=|A|E
\]
且 $A$ 奇异可得 $A^*A=O$.更直接地,$A\boldsymbol{1}=0$ 表明 $A$ 的各列线性关系系数全为 $1$.当 $r(A)=n-1$ 时,任一非零的伴随矩阵行都张成 $\ker A$,故第一行
\[
(A_{11},A_{12},\ldots,A_{1n})
\]
必与 $(1,1,\ldots,1)$ 成比例;当 $r(A)\leq n-2$ 时伴随矩阵为零.两种情形都得到
\[
\boxed{A_{11}=A_{12}=\cdots=A_{1n}}.
\]
\end{solution}


\newpage

\section{对偶空间与商空间}

\begin{definition}
设 $V$ 为数域 $F$ 上的线性空间.若映射 $f:V\to F$ 满足
\[
f(\alpha+\beta)=f(\alpha)+f(\beta),
\qquad
f(k\alpha)=kf(\alpha),
\]
则称 $f$ 为 $V$ 上的线性函数或线性泛函.所有线性函数组成的空间记为
\[
V^*=\operatorname{Hom}(V,F),
\]
称为 $V$ 的对偶空间.
\end{definition}

若 $V$ 为有限维且 $\alpha_1,\ldots,\alpha_n$ 为一组基,则任意 $f\in V^*$ 完全由 $f(\alpha_1),\ldots,f(\alpha_n)$ 决定.因此有:

\begin{proposition}[\markstar]
若 $\dim V<\infty$,则
\[
\dim V^*=\dim V.
\]
\end{proposition}

\begin{definition}
给定 $V$ 的基 $\alpha_1,\ldots,\alpha_n$,若 $V^*$ 的基 $f_1,\ldots,f_n$ 满足
\[
f_i(\alpha_j)=\delta_{ij},
\]
则称 $f_1,\ldots,f_n$ 为关于该基的对偶基.
\end{definition}

\begin{theorem}[\markstar\ 自然映射到二重对偶]
若 $V$ 为有限维线性空间,则映射
\[
\tau:V\to V^{**},
\qquad
\alpha\longmapsto\varphi_\alpha,
\]
其中
\[
\varphi_\alpha(f)=f(\alpha),
\qquad f\in V^*,
\]
是自然同构.
\end{theorem}

%例题4.5.1
\begin{example}
设
\[
\alpha_1=(1,-1,4),\quad
\alpha_2=(1,0,2),\quad
\alpha_3=(0,3,-5)
\]
是 $\mathbb{R}^3$ 的一组基.求其对偶基.
\end{example}
\exampleblank

\begin{solution}
令
\[
P=(\alpha_1,\alpha_2,\alpha_3)
=
\begin{pmatrix}
1&1&0\\
-1&0&3\\
4&2&-5
\end{pmatrix}.
\]
若把线性泛函写成行向量,则对偶基的行向量正是 $P^{-1}$ 的各行.计算得
\[
P^{-1}=
\begin{pmatrix}
-6&5&3\\
7&-5&-3\\
-2&2&1
\end{pmatrix}.
\]
故
\[
\boxed{
\begin{aligned}
f_1(x)&=-6x_1+5x_2+3x_3,\\
f_2(x)&=7x_1-5x_2-3x_3,\\
f_3(x)&=-2x_1+2x_2+x_3.
\end{aligned}}
\]
\end{solution}

%例题4.5.2
\begin{example}
设 $a_1,\ldots,a_n\in\mathbb{R}$ 两两不同,令
\[
f_i(x)=\prod_{j\neq i}(x-a_j),
\qquad i=1,\ldots,n.
\]
\begin{enumerate}
  \item 证明 $f_1,\ldots,f_n$ 是 $\mathbb{R}[x]_n$ 的一组基;
  \item 求由基 $f_1,\ldots,f_n$ 到基 $1,x,\ldots,x^{n-1}$ 的过渡矩阵;
  \item 求 $f_1,\ldots,f_n$ 的对偶基.
\end{enumerate}
\end{example}
\exampleblank

\begin{solution}
若
\[
\displaystyle\sum\limits_{i=1}^n c_if_i(x)=0,
\]
代入 $x=a_k$ 得
\[
c_kf_k(a_k)=0.
\]
因 $a_i$ 两两不同,$f_k(a_k)\neq0$,故所有 $c_k=0$.又共有 $n$ 个向量,故构成 $\mathbb{R}[x]_n$ 的一组基.

设
\[
e_r^{(i)}
=
 e_r(a_1,\ldots,\widehat{a_i},\ldots,a_n)
\]
表示删去 $a_i$ 后的第 $r$ 个初等对称多项式,则
\[
f_i(x)
=
\sum\limits_{k=0}^{n-1}
(-1)^{n-1-k}e_{n-1-k}^{(i)}x^k.
\]
因此过渡矩阵第 $i$ 列为
\[
\left(
(-1)^{n-1}e_{n-1}^{(i)},
(-1)^{n-2}e_{n-2}^{(i)},
\ldots,
-e_1^{(i)},
1
\right)^{\mathsf T}.
\]

对偶基最简洁地写成评价泛函:
\[
\boxed{
\phi_i(p)
=
\dfrac{p(a_i)}{f_i(a_i)}
=
\dfrac{p(a_i)}{\displaystyle\prod_{j\neq i}(a_i-a_j)}
}.
\]
显然 $\phi_i(f_j)=\delta_{ij}$.
\end{solution}

%例题4.5.3
\begin{example}
设 $V$ 为 $n$ 维线性空间.证明:$n$ 个线性函数 $f_1,\ldots,f_n$ 构成 $V^*$ 的一组基,当且仅当不存在非零向量 $x\in V$ 使
\[
f_1(x)=\cdots=f_n(x)=0.
\]
\end{example}
\exampleblank

\begin{solution}
定义线性映射
\[
\Phi:V\to F^n,
\qquad
x\longmapsto(f_1(x),\ldots,f_n(x))^{\mathsf T}.
\]
题中条件正是 $\ker\Phi=\{0\}$.由于 $\dim V=\dim F^n=n$,这等价于 $\Phi$ 可逆.取 $V$ 的任一基后,$\Phi$ 的矩阵各行正是 $f_i$ 的坐标行,因此 $\Phi$ 可逆当且仅当这些行线性无关,即 $f_1,\ldots,f_n$ 构成 $V^*$ 的一组基.
\end{solution}

\begin{definition}
设 $W$ 为线性空间 $V$ 的子空间,$\alpha,\beta\in V$.若
\[
\alpha-\beta\in W,
\]
称 $\alpha,\beta$ 模 $W$ 同余,记为
\[
\alpha\equiv\beta\pmod W.
\]
同余类
\[
\overline\alpha=\alpha+W
\]
组成的集合
\[
V/W=\{\alpha+W\mid\alpha\in V\}
\]
按
\[
(\alpha+W)+(\beta+W)=(\alpha+\beta)+W,
\]
\[
c(\alpha+W)=c\alpha+W
\]
定义运算,称为商空间.
\end{definition}

\begin{theorem}[\markstar]
设 $W$ 为 $V$ 的子空间,则:
\begin{enumerate}
  \item $V/W$ 是数域 $F$ 上的线性空间;
  \item 若 $\varepsilon_1,\ldots,\varepsilon_r$ 为 $W$ 的基,并扩充成 $V$ 的基
  \[
  \varepsilon_1,\ldots,\varepsilon_r,\varepsilon_{r+1},\ldots,\varepsilon_n,
  \]
  则
  \[
  \overline\varepsilon_{r+1},\ldots,\overline\varepsilon_n
  \]
  是 $V/W$ 的一组基,且
  \[
  \dim(V/W)=\dim V-\dim W;
  \]
  \item 自然映射
  \[
  \varphi_W:V\to V/W,
  \qquad
  \alpha\mapsto\alpha+W
  \]
  为满线性映射,且 $\ker\varphi_W=W$.
\end{enumerate}
\end{theorem}

\begin{theorem}[\marktwostar\ 第一同构定理]
设 $\psi:V_1\to V_2$ 为线性映射,则存在自然同构
\[
V_1/\ker\psi\cong\operatorname{Im}\psi,
\]
由
\[
\alpha+\ker\psi\longmapsto\psi(\alpha)
\]
给出.
\end{theorem}

\begin{theorem}[\markstar\ 核的泛性]
设 $\sigma:V_1\to V_2$ 为线性映射,$W\subseteq\ker\sigma$.则存在唯一线性映射
\[
\sigma':V_1/W\to V_2
\]
使
\[
\sigma=\sigma'\circ\varphi_W.
\]
\end{theorem}

%例题4.5.4
\begin{example}
设 $V$ 为数域 $F$ 上的线性空间,$W$ 为其子空间,$V^*$ 为对偶空间,
\[
W^0=\{f\in V^*\mid f(W)=0\}
\]
为 $W$ 的零化子.证明
\[
W^*\cong V^*/W^0.
\]
\end{example}
\exampleblank

\begin{solution}
定义限制映射
\[
R:V^*\to W^*,
\qquad
R(f)=f|_W.
\]
其核显然为 $W^0$.有限维情形下,任意 $W$ 上线性泛函都可把 $W$ 的基扩充为 $V$ 的基后延拓到 $V$,故 $R$ 为满射.由第一同构定理,
\[
\boxed{V^*/W^0\cong W^*}.
\]
\end{solution}

%例题4.5.5
\begin{example}
设
\[
f(x)=x^4+3x^3-x^2-4x-3,
\]
\[
g(x)=3x^3+10x^2+2x-3.
\]
\begin{enumerate}
  \item 求 $(f,g)$;
  \item 令
  \[
  U=\{uf+vg\mid u,v\in F[x]\},
  \]
  求 $F[x]/U$ 的维数.
\end{enumerate}
\end{example}
\exampleblank

\begin{solution}
在通常特征零数域上,
\[
f=(x+3)(x^3-x-1),
\]
\[
g=(x+3)(3x^2+x-1),
\]
而后两个因子互素,因此
\[
\boxed{(f,g)=x+3}.
\]
由 Bézout 定理,
\[
U=(f,g)=(x+3),
\]
故
\[
F[x]/U\cong F[x]/(x+3)
\]
为一维空间.因此
\[
\boxed{\dim(F[x]/U)=1}.
\]
\end{solution}

%例题4.5.6
\begin{example}
设 $A,B$ 是有限维向量空间 $V$ 的两个子空间.
\begin{enumerate}
  \item 证明
  \[
  A/(A\cap B)\cong(A+B)/B;
  \]
  \item 证明自然同构
  \[
  (A+B)/(A\cap B)
  \cong
  A/(A\cap B)\oplus B/(A\cap B).
  \]
\end{enumerate}
\end{example}
\exampleblank

\begin{solution}
定义
\[
\phi:A\to(A+B)/B,
\qquad
\phi(a)=a+B.
\]
其核为 $A\cap B$,像为整个 $(A+B)/B$,由第一同构定理得
\[
\boxed{A/(A\cap B)\cong(A+B)/B}.
\]

再定义
\[
\Psi:
A/(A\cap B)\oplus B/(A\cap B)
\to
(A+B)/(A\cap B)
\]
为
\[
\Psi\left(a+(A\cap B),b+(A\cap B)\right)
=a+b+(A\cap B).
\]
它良定义且显然满射.若其像为零,则 $a+b\in A\cap B$.因为 $a\in A,b\in B$,由此推出 $a,b\in A\cap B$,故核为零.因此 $\Psi$ 为同构.
\end{solution}
